\documentclass[10pt,a4paper,sans]{moderncv}
\moderncvstyle{classic}
\moderncvcolor{blue}
\usepackage[utf8]{inputenc}
\usepackage[scale=0.85]{geometry}
\usepackage{helvet}

\name{Alexandro}{Disla}
\title{Suivi-Évaluation et Redevabilité (MEAL)}
\address{Rue Saint-Louis Jeanty \#14, Route de Frère}{Port-au-Prince, Haïti}{}
\phone[mobile]{(+509) 4148-3700}
\email{alexandrodisla@hotmail.com}
\extrainfo{GitHub: \url{https://github.com/AD0791}}

\begin{document}
\makecvtitle

\section{Profil Professionnel}
\cvitem{}{Professionnel MEAL avec \textbf{cinq ans d'expérience en ONG en Haïti}, sur trois secteurs --- santé et VIH à la Caris Foundation International, eau et assainissement chez HANWASH, éducation chez Anseye Pou Ayiti --- précédés de huit ans de suivi d'indicateurs publics au Ministère de la Planification. Mon travail porte sur la chaîne qui relie une information recueillie à une décision : un enregistrement juste dès la saisie, un suivi qui ne s'arrête qu'à la clôture, des bases tenues à jour, une analyse de tendances et des rapports périodiques qui résistent à la vérification. Trois ans sur les données d'un programme VIH m'ont appris à traiter l'information sensible avec la discrétion qu'elle exige. \textbf{Créole haïtien et français maternels}, anglais courant à l'écrit comme à l'oral. Résidant à Route de Frère, à proximité de Pétion-Ville ; disponible pour les visites de terrain et le travail le samedi.}

\section{Compétences Clés}
\cvitem{}{
\begin{itemize}
    \item \textbf{Enregistrement, suivi et clôture :} Tenue de registres structurés et de flux standardisés entre sites ; suivi des actions correctives jusqu'à leur clôture effective, et non jusqu'à leur simple consignation ; coordination avec les équipes de programme, les agents de terrain et les partenaires gouvernementaux (MSPP, DINEPA, MEF).
    \item \textbf{Informations sensibles \& confidentialité :} Données d'un programme de santé et de VIH ; contrôle d'accès par rôle et stockage sécurisé dans les systèmes que je conçois ; dédoublonnage et correction contrôlée des registres de bénéficiaires.
    \item \textbf{Bases de données, qualité \& analyse :} Protocoles d'assurance qualité des données (DQA) établis et supervisés, vérification par remontée aux registres sources ; administration de bases relationnelles ; catégorisation stable des données et analyse de tendances ; \textbf{Power BI} et Excel avancé.
    \item \textbf{Rapportage :} Rapports d'activité automatisés en Python et SQL (temps de traitement manuel réduit de 40 \%) ; tableaux de bord de suivi hebdomadaire ; production et soumission d'indicateurs destinés à l'USAID via la plateforme DATIM ; guides techniques et notes de restitution pour des équipes non techniques.
    \item \textbf{Terrain \& communication :} Direction technique d'agents de collecte sur des sites multiples ; formation d'agents de terrain et d'équipes aux outils et procédures ; restitution en créole, en français et en anglais.
\end{itemize}}

\section{Expériences en Organisations Non Gouvernementales}

\cventry{Oct 2021 -- Août 2024}{Officier de Suivi \& Évaluation}{Impact Youth Project, Caris Foundation International}{Haïti}{}{
\begin{itemize}
    \item \textbf{Suivi multi-sites :} Conduite du cycle de suivi d'un programme de santé et de VIH sur plusieurs sites, avec standardisation des formulaires et des flux de données entre sites et consolidation des remontées dans les délais.
    \item \textbf{Données sensibles :} Collecte numérique (CommCare, HIVHaiti) et bases MySQL intégrées portant sur des bénéficiaires d'un programme VIH, sous les exigences de confidentialité propres à ce type de données.
    \item \textbf{Assurance qualité (DQA) :} Établissement et supervision des protocoles de Data Quality Assessment, avec remontée des chiffres rapportés jusqu'aux registres sources et suivi des actions correctives jusqu'à clôture.
    \item \textbf{Encadrement technique :} Direction technique des agents de collecte sur l'ensemble des sites --- répartition du travail, contrôle qualité, formation aux procédures --- et référent technique du personnel de suivi-évaluation.
    \item \textbf{Rapportage :} Production et soumission des indicateurs MER destinés à l'USAID via DATIM, la plateforme du PEPFAR bâtie sur DHIS2 ; rapports d'activité automatisés en Python et SQL, réduisant de 40 \% le temps de traitement manuel ; tableaux de bord Power BI et Quarto pour le suivi hebdomadaire.
\end{itemize}}

\cventry{Oct 2023 -- Jan 2026}{Consultant MEAL \& Analyste de Données}{HANWASH}{Haïti / Télétravail}{}{
\begin{itemize}
    \item Raffinement des plans MEAL et des tableaux de suivi des indicateurs alignés sur les normes JMP (OMS/UNICEF) pour des programmes d'eau, d'hygiène et d'assainissement.
    \item Déploiement d'enquêtes géospatiales sur mWater, administration de la base et résolution des incidents remontés du terrain pour préserver l'intégrité des séries.
    \item Tableaux de bord automatisés (mWater, Power BI) connectés aux API de collecte ; rédaction des guides techniques et formation des équipes à leur usage autonome.
\end{itemize}}

\cventry{Jan 2026 -- Mars 2026}{Analyste de Données \& Consultant MEAL}{Anseye Pou Ayiti}{Haïti}{}{
\begin{itemize}
    \item Définition de la stratégie de collecte et élaboration des indicateurs de performance de trois programmes éducatifs ; codage des instruments XLSForm pour ODK avec logiques de saut et contraintes de validation.
    \item Registre de bénéficiaires alimenté par ODK : recherche, correction contrôlée des enregistrements, détection des doublons et export Excel ; formation des agents de terrain aux outils de collecte.
\end{itemize}}

\section{Expériences dans le Secteur Public et en Systèmes d'Information}

\cventry{Nov 2015 -- Août 2023}{Économiste --- Direction de Planification Économique et Sociale}{Ministère de la Planification et de la Coopération Externe (MPCE)}{Port-au-Prince}{}{Rapports de suivi de l'exécution du Plan Triennal d'Investissement 2014--2016 : indicateurs de projets publics sur l'ensemble des secteurs et des départements ; unité statistique et modélisation.}

\cventry{Févr 2026 -- Présent}{Ingénieur Logiciel (Fullstack)}{Tekkod LLC}{Télétravail}{}{Application mobile fiable hors ligne et sécurisée pour un usage terrain ; tableaux de bord Looker Studio. Contrat en phase de clôture.}

\cventry{Mars 2021 -- Mai 2024}{Développeur Backend \& Ingénieur de Données}{Tekkod LLC}{Télétravail}{}{Administration de bases relationnelles en accès concurrent ; pipelines de données ; services sécurisés avec contrôle d'accès par rôle.}

\cventry{Nov 2015 -- Mars 2021}{Consultant Logiciel \& Intégration de Données}{CassionSoft (Projet UNOPS)}{Haïti}{}{Intégration de données sur un projet sous contrat des Nations Unies, avec exposition aux procédures et standards documentaires onusiens.}

\section{Formation}
\cvitem{2010--2014}{\textbf{Études Supérieures en Économie Appliquée} --- C.T.P.E.A (Centre de Techniques de Planification et d'Économie Appliquée), Port-au-Prince : cycle de quatre ans en économie et statistique appliquées. \textit{Scolarité complétée et ensemble des examens réussis ; Diplôme d'Études Supérieures en attente de la soutenance du mémoire de sortie --- attestation d'études jointe au dossier.}}
\cvitem{2009--2010}{Baccalauréat (Mention) --- Institution Saint-Louis de Gonzague, Port-au-Prince}

\section{Compétences Techniques}
\cvitem{Microsoft Office}{\textbf{Word, Excel (avancé), PowerPoint, Teams}}
\cvitem{Analyse \& viz}{\textbf{Power BI}, Looker Studio, SQL, Python (Pandas), R (RMarkdown/Quarto)}
\cvitem{Collecte}{\textbf{KoboToolbox, ODK, XLSForm}, CommCare, mWater, DATIM \textit{(soumission de données)}}
\cvitem{Bases de données}{MySQL, PostgreSQL, MongoDB --- conception, administration, sauvegarde, contrôle d'accès}
\cvitem{Langues}{\textbf{Créole haïtien (maternel)}, \textbf{Français (maternel)}, \textbf{Anglais (courant, écrit et oral)}}

\end{document}
